\documentclass[11pt, letterpaper]{article}

% --- Idioma y codificación ---
\usepackage[spanish]{babel}
\usepackage[utf8]{inputenc}
\usepackage[T1]{fontenc}
\usepackage{lmodern}

% --- Matemáticas ---
\usepackage{amsmath, amssymb, amsthm}

% --- Formato ---
\usepackage[margin=2.5cm]{geometry}
\usepackage{parskip}
\usepackage{booktabs}
\usepackage{hyperref}
\hypersetup{colorlinks=true, linkcolor=blue!70!black, urlcolor=blue!70!black}

% --- Entornos de teoremas ---
\theoremstyle{definition}
\newtheorem{definition}{Definición}[section]
\theoremstyle{plain}
\newtheorem{theorem}{Teorema}[section]
\newtheorem{lemma}[theorem]{Lema}
\newtheorem{corollary}[theorem]{Corolario}
\theoremstyle{remark}
\newtheorem{remark}{Observación}[section]

\title{%
  T02 --- Demostración formal:\\[4pt]
  Fórmulas de Knuth para sonda lineal
}
\author{Curso Linux--C--Rust --- B15 Estructuras de datos en C}
\date{}

\begin{document}
\maketitle
\tableofcontents

% ===================================================================
\section{Modelo}
% ===================================================================

\begin{definition}
Array $T[0..m{-}1]$ con $m$ slots. Cada slot está vacío u ocupado por
exactamente un elemento. Para insertar o buscar la clave~$k$, se
examina la secuencia:
\[
  h(k),\; h(k){+}1,\; h(k){+}2,\; \ldots \pmod{m}
\]
hasta encontrar un slot vacío (búsqueda fallida / inserción) o la
clave buscada (búsqueda exitosa). Bajo hashing uniforme,
$h(k) \sim \mathrm{Uniforme}\{0, \ldots, m{-}1\}$.
Notación: $n$ elementos, $m$ capacidad, $\alpha = n/m < 1$.
Contamos \textbf{sondeos} (slots examinados).
\end{definition}

% ===================================================================
\section{Demostración 1: búsqueda fallida
  \texorpdfstring{$\frac{1}{2}\!\left(1 + \frac{1}{(1-\alpha)^2}\right)$}%
  {1/2(1 + 1/(1-alpha)\^2)}}
% ===================================================================

\begin{theorem}\label{thm:miss}
El número esperado de sondeos en una búsqueda fallida con sonda lineal
es:
\[
  E[\text{sondeos fallida}]
  \approx \frac{1}{2}\left(1 + \frac{1}{(1-\alpha)^2}\right)
\]
\end{theorem}

\begin{proof}[Derivación]
Definimos $A_k$ = evento de que se necesitan al menos $k$ sondeos.
Esto ocurre si los $k{-}1$ primeros slots examinados están todos
ocupados. El número esperado de sondeos es:
\[
  E[\text{sondeos}]
  = \sum_{k=1}^{m} \Pr[A_k]
  = 1 + \sum_{k=2}^{m} \Pr[A_k]
\]

Bajo hashing uniforme, la probabilidad de que $k{-}1$ slots
consecutivos específicos estén todos ocupados es:
\[
  \Pr[A_k]
  = \frac{n(n{-}1)\cdots(n{-}k{+}2)}{m(m{-}1)\cdots(m{-}k{+}2)}
  \approx \alpha^{k-1}
\]

La aproximación $\alpha^{k-1}$ ignora las \textbf{correlaciones del
clustering primario}. Un cluster de longitud~$L$ absorbe inserciones
destinadas a $L + 1$ posiciones (las $L$ del cluster más la
inmediatamente posterior), creando un sesgo que hace que los clusters
largos sean más probables.

Knuth demostró, mediante un análisis combinatorio fino (TAOCP Vol.~3,
\S6.4) usando funciones generatrices para la distribución exacta de
longitudes de cluster, que el efecto del clustering transforma
$1/(1{-}\alpha)$ (caso sin clustering) en $1/(1{-}\alpha)^2$:
\[
  E[\text{sondeos fallida}]
  \approx \frac{1}{2}\left(1 + \frac{1}{(1-\alpha)^2}\right) \qedhere
\]
\end{proof}

% ===================================================================
\section{Demostración 2: búsqueda exitosa
  \texorpdfstring{$\frac{1}{2}\!\left(1 + \frac{1}{1-\alpha}\right)$}%
  {1/2(1 + 1/(1-alpha))}}
% ===================================================================

\begin{theorem}\label{thm:hit}
El número esperado de sondeos en una búsqueda exitosa con sonda lineal
es:
\[
  E[\text{sondeos exitosa}]
  \approx \frac{1}{2}\left(1 + \frac{1}{1-\alpha}\right)
\]
\end{theorem}

\begin{proof}
El costo de encontrar un elemento es igual al costo que tuvo
\textbf{insertarlo}. Cuando se insertó el $i$-ésimo elemento, la
tabla contenía $i - 1$ elementos con factor de carga
$\alpha_i = (i{-}1)/m$. El costo de esa inserción (= búsqueda
fallida con $i{-}1$ elementos) fue:
\[
  C_i \approx \frac{1}{2}\left(1 + \frac{1}{(1 - (i{-}1)/m)^2}\right)
\]

El costo promedio de búsqueda exitosa es:
\[
  E[\text{exitosa}]
  = \frac{1}{n} \sum_{i=1}^{n} C_i
  = \frac{1}{2} + \frac{1}{2n}
    \sum_{i=0}^{n-1} \frac{1}{(1 - i/m)^2}
\]

Aproximamos la suma por una integral. Con $x = i/m$, $dx = 1/m$:
\[
  \frac{1}{n} \sum_{i=0}^{n-1} \frac{1}{(1 - i/m)^2}
  \approx \frac{m}{n} \int_0^{n/m} \frac{dx}{(1-x)^2}
  = \frac{1}{\alpha} \int_0^{\alpha} \frac{dx}{(1-x)^2}
\]

Calculamos la integral:
\[
  \int_0^{\alpha} \frac{dx}{(1-x)^2}
  = \left[\frac{1}{1-x}\right]_0^{\alpha}
  = \frac{1}{1-\alpha} - 1
  = \frac{\alpha}{1-\alpha}
\]

Sustituyendo:
\[
  E[\text{exitosa}]
  \approx \frac{1}{2}\left(1 + \frac{1}{\alpha}
    \cdot \frac{\alpha}{1-\alpha}\right)
  = \frac{1}{2}\left(1 + \frac{1}{1-\alpha}\right) \qedhere
\]
\end{proof}

% ===================================================================
\section{Demostración 3: divergencia cuando
  \texorpdfstring{$\alpha \to 1$}{alpha -> 1}}
% ===================================================================

\begin{theorem}\label{thm:divergence}
Ambas fórmulas divergen cuando $\alpha \to 1^-$.
\end{theorem}

\begin{proof}
Inmediata de las fórmulas:

\emph{Búsqueda fallida:}
\[
  \frac{1}{2}\left(1 + \frac{1}{(1-\alpha)^2}\right)
  \xrightarrow{\alpha \to 1^-} +\infty
  \qquad \text{porque } (1{-}\alpha)^2 \to 0^+
\]

\emph{Búsqueda exitosa:}
\[
  \frac{1}{2}\left(1 + \frac{1}{1-\alpha}\right)
  \xrightarrow{\alpha \to 1^-} +\infty
  \qquad \text{porque } (1{-}\alpha) \to 0^+
\]

\textbf{Tasa de divergencia.} La búsqueda fallida diverge
\textbf{mucho más rápido}:
\[
  \frac{E[\text{fallida}]}{E[\text{exitosa}]}
  \approx \frac{1/(1{-}\alpha)^2}{1/(1{-}\alpha)}
  = \frac{1}{1-\alpha}
  \xrightarrow{\alpha \to 1^-} +\infty
\]
Para $\alpha = 0.95$: fallida $\approx 200.5$, exitosa $\approx 10.5$.
Ratio $\approx 19$.
\end{proof}

\begin{remark}
Una tabla casi llena con sonda lineal es catastrófica para búsquedas
de claves inexistentes. Esto motiva que el umbral de
redimensionamiento sea $\alpha \leq 0.7$--$0.75$.
\end{remark}

% ===================================================================
\section{Demostración 4: comparación con hashing uniforme}
% ===================================================================

\begin{theorem}\label{thm:uniform}
Bajo hashing uniforme (sin clustering), los costos son:
\[
  E[\text{fallida, uniforme}] = \frac{1}{1-\alpha}, \qquad
  E[\text{exitosa, uniforme}] = \frac{1}{\alpha}
    \ln\frac{1}{1-\alpha}
\]
que son estrictamente menores que los de sonda lineal.
\end{theorem}

\begin{proof}
\emph{Caso fallido.} La probabilidad de que el primer sondeo caiga en
un slot ocupado es~$\alpha$. Si es así, el segundo sondeo (aleatorio
independiente) cae en un slot ocupado con probabilidad
$(n{-}1)/(m{-}1) \approx \alpha$. Se necesitan exactamente $k$
sondeos con probabilidad $\alpha^{k-1}(1{-}\alpha)$ (distribución
geométrica):
\[
  E[\text{sondeos}] = \frac{1}{1-\alpha}
\]

\emph{Caso exitoso.} Análogo a la Demostración~2, pero con costo de
inserción $1/(1{-}\alpha_i)$ en vez de $1/(1{-}\alpha_i)^2$:
\[
  E[\text{exitosa}]
  = \frac{1}{n}\sum_{i=0}^{n-1}\frac{1}{1{-}i/m}
  \approx \frac{1}{\alpha}\int_0^{\alpha}\frac{dx}{1{-}x}
  = \frac{1}{\alpha}\ln\frac{1}{1{-}\alpha} \qedhere
\]
\end{proof}

El clustering primario cuesta un factor de $1/(1{-}\alpha)$ adicional:

\begin{table}[h]
\centering
\caption{Costo del clustering: sonda lineal vs hashing uniforme.}
\begin{tabular}{cccc}
\toprule
$\alpha$ & Fallida (lineal) & Fallida (uniforme)
  & Ratio \\
\midrule
0.50 & 2.50  & 2.00 & 1.25 \\
0.70 & 6.06  & 3.33 & 1.82 \\
0.80 & 13.0  & 5.00 & 2.60 \\
0.90 & 50.5  & 10.0 & 5.05 \\
0.95 & 200.5 & 20.0 & 10.0 \\
\bottomrule
\end{tabular}
\end{table}

% ===================================================================
\section{Demostración 5: peor caso
  \texorpdfstring{$\Theta(n^2)$}{Theta(n\^2)}}
% ===================================================================

\begin{theorem}\label{thm:worst}
En el peor caso, $n$ inserciones con sonda lineal cuestan
$\Theta(n^2)$ sondeos totales.
\end{theorem}

\begin{proof}
Considérese el caso en que todas las $n$ claves hashean al mismo
slot~$j$. La $i$-ésima inserción requiere $i$ sondeos (recorrer los
$i{-}1$ slots ocupados más el slot vacío):
\[
  \text{Sondeos totales}
  = \sum_{i=1}^{n} i
  = \frac{n(n+1)}{2}
  = \Theta(n^2)
\]

Cada búsqueda de la clave insertada en la posición~$i$ cuesta $i$
sondeos. La búsqueda exitosa promedio:
\[
  E[\text{exitosa}]
  = \frac{1}{n}\sum_{i=1}^{n} i
  = \frac{n+1}{2}
  = \Theta(n) \qedhere
\]
\end{proof}

\begin{remark}
Bajo SUHA, el costo esperado es $O(1/(1{-}\alpha))$. El gap entre
caso esperado $O(1)$ (con $\alpha$ constante) y peor caso $\Theta(n)$
subraya la importancia de una buena función hash.
\end{remark}

% ===================================================================
\section{Verificación numérica}
% ===================================================================

\begin{table}[h]
\centering
\begin{tabular}{cccc}
\toprule
$\alpha$
  & $\frac{1}{2}(1 + 1/(1{-}\alpha))$
  & $\frac{1}{2}(1 + 1/(1{-}\alpha)^2)$
  & Ratio \\
\midrule
0.10 & 1.056  & 1.117   & 1.06 \\
0.25 & 1.167  & 1.389   & 1.19 \\
0.50 & 1.500  & 2.500   & 1.67 \\
0.70 & 2.167  & 6.056   & 2.80 \\
0.75 & 2.500  & 8.500   & 3.40 \\
0.80 & 3.000  & 13.00   & 4.33 \\
0.90 & 5.500  & 50.50   & 9.18 \\
0.95 & 10.50  & 200.5   & 19.1 \\
\bottomrule
\end{tabular}
\end{table}

% ===================================================================
\section{Resumen de resultados}
% ===================================================================

\begin{table}[h]
\centering
\begin{tabular}{llp{5.5cm}}
\toprule
\textbf{Resultado} & \textbf{Técnica} & \textbf{Clave} \\
\midrule
Fallida: $\frac{1}{2}(1 + 1/(1{-}\alpha)^2)$
  & Análisis de Knuth
  & Clustering eleva $1/(1{-}\alpha)$ a $1/(1{-}\alpha)^2$ \\
Exitosa: $\frac{1}{2}(1 + 1/(1{-}\alpha))$
  & Promedio de costos de inserción
  & $\int_0^\alpha dx/(1{-}x)^2 = \alpha/(1{-}\alpha)$ \\
Divergencia en $\alpha \to 1$
  & Análisis de límites
  & Fallida diverge como $1/(1{-}\alpha)^2$ \\
Clustering: ${\times}\,1/(1{-}\alpha)$ vs uniforme
  & Comparación de modelos
  & Sonda lineal vs hashing uniforme \\
Peor caso $\Theta(n^2)$
  & Todas las claves al mismo slot
  & $\sum i = n(n{+}1)/2$ \\
\bottomrule
\end{tabular}
\end{table}

\end{document}
